\chapter{Introduction}
\label{sec:introduction}

Causal inference is a field of study that aims to understand and estimate the 
causal effects of intervetions. Many fundamental questions in psychology, epidemiology, 
economics and political science are causal in nature. Questions such as \emph{Will 
increasing the minimum wage lead to higher employment?} \cite{cardMinimumWagesEmployment1993} or \emph{
    Does initiating hormonal therapy increase the risk of coronary heart disease (CHD) in postmenopausal women?
} \cite{hernanObservationalStudiesAnalyzed2008}. The common thread in these questions is that they 
are querying about an intervention and its effect on an outcome. Therefore, to answer them, it is not enough 
to simply find an association between those intervetions and the outcomes, but we need to know or assume 
something about
the \emph{structure} of the data generating process or the causal diagram 
and using that structure to identify the 
causal target of interest. 

Knowing the structure of the data generating process 
is not an easy task. 
It often requires an expert knowledge of the domain and assumptions about how the 
world works. One way to mitigate 
the problem of not knowing the structure of the data generating process is to conduct a randomized controlled trial (RCT). 
RCTs are often considered the gold standard for causal inference 
because they allow researchers to control for confounding variables and establish a 
clear cause-and-effect relationship between the treatment and the outcome. 
However, in many cases,
conducting randomized experiments is often infeasible, unethical or expensive and therefore, 
researchers often have to rely on observational data to estimate 
causal effects. 

Causal inference from observational datasets usually involves two 
main levels of  complexity.
%There are two main levels of complexity when it comes to 
%conducting causal inference with observational data. 
The first level is when the 
cuasal diagram is unknown and we wish to discover it from the data. This is the topic of 
causal discovery, an active area of research and not the subject of this thesis. 
The second level of complexity is when the causal diagram is known but we have unmeasured confounders, i.e. 
there are variables that affect both the treatment and the outcome but are not observed in the data. This is the 
focus of this thesis. Furthermore, we are going to limit our focus to longitudinal data, a common research design where 
the same individuals are observed at multiple time points. In many contemporary applications, 
these repeated measurements are collected at high frequency in daily life through experience sampling, 
smartphone-based assessments, and wearable sensors. The growing availability of such intensive 
longitudinal data has made methods for longitudinal causal inference increasingly important \cite{myin-germeysExperienceSamplingMethodology2018}.
 
Estimating causal effects in such designs where 
treatments and confounders can vary over time is particularly challenging 
because of the presence of treatment-confounder feedback, i.e.
the treatments at previous time points both affect the confounders at the next time points 
and the outcomes. To further define the study's boundaries, it is worth mentioning 
that this paper focuses on \emph{discrete} time periods\cite{gillCausalInferenceComplex2001} and assumes a lack of anticipation effects, 
where treatments at future points do not impact potential outcomes at the current time point \cite{leeFindingDynamicTreatment2012}.

\section{Research question}
This thesis aims to compare the performance and assumptions of two methods that have been 
proposed to estimate the causal effects of time-varying treatments 
in the presence of \emph{time-invariant} unmeasured confounders. The first method is the FE-MSM method 
proposed by \cite{blackwellEffectPoliticalAdvertising2024} and the second method is the 
SeqGPLVM method proposed by \cite{hattSequentialDeconfoundingCausal2024}. 
Both approaches are comparable as they aim to estimate the 
causal effects of time-varying treatments in the presence of time-invariant 
unmeasured confounders espechially when previous treatments influence time-varying confounders. 
Both methods are applicable to similar setting with  
a single treatment available at each time point and a single outcome at the end of the study.
Moreover, they both invlove a unit-specific unobserved heterogeneity component that 
is supposed to capture the effect of the
unmeasured confounders. 

They diverge, however, in their characterization of this 
unobserved heterogeneity component, the assumptions they make about it and the scope of 
applicability of their methods. 
Whereas FE-MSM—rooted in econometrics and political science literature—conceptualizes this 
component as a frequentist fixed effect and a nuisance parameter, 
the machine-learning-based SeqGPLVM views it 
as a latent variable, a common framework in machine learning and psychometrics literature. Therefore, putting these two methods 
together in the same paper using a unified notation and vocabulary is an effort to make a 
bridge between different fields of research and to highlight the 
assumptions and implications of these two methods in a same framework. 

\section{Contribution}
The main contribution of this thesis is as follows\footnote{Code available at \url{https://github.com/Ali-Setareh/SeqGPLVM_FEMSM}}:
\begin{itemize}
    \item To place the FE-MSM and SeqGPLVM approaches in a common notation and conceptual framework, making their assumptions, targets, and implications more directly comparable.
    
    \item To evaluate FE-MSM and SeqGPLVM on the same synthetic data-generating processes and to compare their performance under a common simulation design.
    
    \item To make the identifying and modeling assumptions underlying both methods more explicit.
    \item To provide reproducible and publicly available code for both methods. In particular, this thesis delivers, to the best of my knowledge, the first modular GPyTorch implementation of SeqGPLVM, designed to accommodate different treatment types and intended for practical reuse beyond the original paper.
    
    \item To relate the sequential deconfounding framework to the deconfounder literature, in particular to the \emph{blessings of multiple causes} idea of Wang and Blei, and to discuss how critiques such as those of D'Amour extend to the time-varying treatment setting.
\end{itemize}

\section{Thesis structure}
The rest of the thesis is structured as follows. Chapter ~\ref{chap:literature_review} reviews the 
relevant literature on longitudinal causal inference with unmeasured confounders.,
Chapter ~\ref{chap:method}, describes the FE-MSM method and the 
SeqGPLVM method in detail, including their assumptions and estimation procedures. 
Chapter ~\ref{chap:simulation_study}, presents a simulation study comparing the performance of these 
two methods under different scenarios. Finally, chapter ~\ref{chap:discussion}, discusses the 
implications of findings and suggest directions for future research.

\section{contribution}
\begin{itemize}
    \item I test both the FE-MSM and SeqGPLVM methods on a same 
    synthetic data and compare their performance. 
    \item by presenting these two methods in a same paper, 
    I can also highlight the differences and similarities between them 
    and perhaps provide some insights about when one method might be preferred over the other. and also 
    encourage future research to explore the connections between methods developed in different fields.
    \item I also make the assumptions behind these two methods more explicit and perhaps discuss 
    about the plausibility of these assumptions in real world settings. 
    \item I provide replicable code for both methods and make it available to the public. 
    \item connection with the blessing of multiple causes paper by Wang and Blei. specifically, 
    the analogous assumption of no single cause confounding in the time varying treatment setting 
    and also the critique made by D'amor about the deconfounder paper 
    and how it applies to the time varying treatment setting. 
    \item related to the previous point, in the paper towards clarifying the
    theory of deconfounder by wang and Blei, they accepted that identifying the potential 
    outcomes for the complete sets of treatments is not possible but it is 
    possible to identify the potential outcomes for subsets of treatments. 
    This is similar to the fe-msm approach where we only identify the potential outcomes for subsets of treatments 
    but they did not explicitly discuss this point in their paper. I can make this connection more explicit. 
\end{itemize}

    